% Options for packages loaded elsewhere
\PassOptionsToPackage{unicode}{hyperref}
\PassOptionsToPackage{hyphens}{url}
\documentclass[
  a4paper,
]{ltjsarticle}
\usepackage{xcolor}
\usepackage[margin=25mm]{geometry}
\usepackage{amsmath,amssymb}
\setcounter{secnumdepth}{-\maxdimen} % remove section numbering
\usepackage{iftex}
\ifPDFTeX
  \usepackage[T1]{fontenc}
  \usepackage[utf8]{inputenc}
  \usepackage{textcomp} % provide euro and other symbols
\else % if luatex or xetex
  \usepackage{unicode-math} % this also loads fontspec
  \defaultfontfeatures{Scale=MatchLowercase}
  \defaultfontfeatures[\rmfamily]{Ligatures=TeX,Scale=1}
\fi
\usepackage{lmodern}
\ifPDFTeX\else
  % xetex/luatex font selection
\fi
% Use upquote if available, for straight quotes in verbatim environments
\IfFileExists{upquote.sty}{\usepackage{upquote}}{}
\IfFileExists{microtype.sty}{% use microtype if available
  \usepackage[]{microtype}
  \UseMicrotypeSet[protrusion]{basicmath} % disable protrusion for tt fonts
}{}
\makeatletter
\@ifundefined{KOMAClassName}{% if non-KOMA class
  \IfFileExists{parskip.sty}{%
    \usepackage{parskip}
  }{% else
    \setlength{\parindent}{0pt}
    \setlength{\parskip}{6pt plus 2pt minus 1pt}}
}{% if KOMA class
  \KOMAoptions{parskip=half}}
\makeatother
\usepackage{color}
\usepackage{fancyvrb}
\newcommand{\VerbBar}{|}
\newcommand{\VERB}{\Verb[commandchars=\\\{\}]}
\DefineVerbatimEnvironment{Highlighting}{Verbatim}{commandchars=\\\{\}}
% Add ',fontsize=\small' for more characters per line
\newenvironment{Shaded}{}{}
\newcommand{\AlertTok}[1]{\textcolor[rgb]{1.00,0.00,0.00}{\textbf{#1}}}
\newcommand{\AnnotationTok}[1]{\textcolor[rgb]{0.38,0.63,0.69}{\textbf{\textit{#1}}}}
\newcommand{\AttributeTok}[1]{\textcolor[rgb]{0.49,0.56,0.16}{#1}}
\newcommand{\BaseNTok}[1]{\textcolor[rgb]{0.25,0.63,0.44}{#1}}
\newcommand{\BuiltInTok}[1]{\textcolor[rgb]{0.00,0.50,0.00}{#1}}
\newcommand{\CharTok}[1]{\textcolor[rgb]{0.25,0.44,0.63}{#1}}
\newcommand{\CommentTok}[1]{\textcolor[rgb]{0.38,0.63,0.69}{\textit{#1}}}
\newcommand{\CommentVarTok}[1]{\textcolor[rgb]{0.38,0.63,0.69}{\textbf{\textit{#1}}}}
\newcommand{\ConstantTok}[1]{\textcolor[rgb]{0.53,0.00,0.00}{#1}}
\newcommand{\ControlFlowTok}[1]{\textcolor[rgb]{0.00,0.44,0.13}{\textbf{#1}}}
\newcommand{\DataTypeTok}[1]{\textcolor[rgb]{0.56,0.13,0.00}{#1}}
\newcommand{\DecValTok}[1]{\textcolor[rgb]{0.25,0.63,0.44}{#1}}
\newcommand{\DocumentationTok}[1]{\textcolor[rgb]{0.73,0.13,0.13}{\textit{#1}}}
\newcommand{\ErrorTok}[1]{\textcolor[rgb]{1.00,0.00,0.00}{\textbf{#1}}}
\newcommand{\ExtensionTok}[1]{#1}
\newcommand{\FloatTok}[1]{\textcolor[rgb]{0.25,0.63,0.44}{#1}}
\newcommand{\FunctionTok}[1]{\textcolor[rgb]{0.02,0.16,0.49}{#1}}
\newcommand{\ImportTok}[1]{\textcolor[rgb]{0.00,0.50,0.00}{\textbf{#1}}}
\newcommand{\InformationTok}[1]{\textcolor[rgb]{0.38,0.63,0.69}{\textbf{\textit{#1}}}}
\newcommand{\KeywordTok}[1]{\textcolor[rgb]{0.00,0.44,0.13}{\textbf{#1}}}
\newcommand{\NormalTok}[1]{#1}
\newcommand{\OperatorTok}[1]{\textcolor[rgb]{0.40,0.40,0.40}{#1}}
\newcommand{\OtherTok}[1]{\textcolor[rgb]{0.00,0.44,0.13}{#1}}
\newcommand{\PreprocessorTok}[1]{\textcolor[rgb]{0.74,0.48,0.00}{#1}}
\newcommand{\RegionMarkerTok}[1]{#1}
\newcommand{\SpecialCharTok}[1]{\textcolor[rgb]{0.25,0.44,0.63}{#1}}
\newcommand{\SpecialStringTok}[1]{\textcolor[rgb]{0.73,0.40,0.53}{#1}}
\newcommand{\StringTok}[1]{\textcolor[rgb]{0.25,0.44,0.63}{#1}}
\newcommand{\VariableTok}[1]{\textcolor[rgb]{0.10,0.09,0.49}{#1}}
\newcommand{\VerbatimStringTok}[1]{\textcolor[rgb]{0.25,0.44,0.63}{#1}}
\newcommand{\WarningTok}[1]{\textcolor[rgb]{0.38,0.63,0.69}{\textbf{\textit{#1}}}}
\usepackage{longtable,booktabs,array}
\newcounter{none} % for unnumbered tables
\usepackage{calc} % for calculating minipage widths
% Correct order of tables after \paragraph or \subparagraph
\usepackage{etoolbox}
\makeatletter
\patchcmd\longtable{\par}{\if@noskipsec\mbox{}\fi\par}{}{}
\makeatother
% Allow footnotes in longtable head/foot
\IfFileExists{footnotehyper.sty}{\usepackage{footnotehyper}}{\usepackage{footnote}}
\makesavenoteenv{longtable}
\setlength{\emergencystretch}{3em} % prevent overfull lines
\providecommand{\tightlist}{%
  \setlength{\itemsep}{0pt}\setlength{\parskip}{0pt}}
\usepackage{bookmark}
\IfFileExists{xurl.sty}{\usepackage{xurl}}{} % add URL line breaks if available
\urlstyle{same}
\hypersetup{
  pdftitle={Chapter 3: Complex-Plane Readout Figures --- Three Grids: Step 0{,} Final Step{,} and Zoom into the Condensed Center},
  pdfauthor={Noriaki Kihara (WF System Co., Ltd.) ORCID 0009-0004-6753-4020},
  hidelinks,
  pdfcreator={LaTeX via pandoc}}

\title{Chapter 3: Complex-Plane Readout Figures --- Three Grids: Step 0,
Final Step, and Zoom into the Condensed Center}
\author{Noriaki Kihara (WF System Co., Ltd.) ORCID 0009-0004-6753-4020}
\date{September 5, 2026 Version DOI 10.5281/zenodo.22317636}

\begin{document}
\maketitle

(Chapter 3 of the reproduction papers for the N=3..40 stage-1+2+3 sweep.
Shares the Concept DOI 10.5281/zenodo.22317635 with the overview paper;
Version DOI 10.5281/zenodo.22317636. Equation numbers continue from
Chapters 1--2 (Eqs. 1--25). Code blocks are quoted verbatim; comments
inside them remain in Japanese, as in the originals.)

\subsection{1. Purpose}\label{purpose}

This chapter describes the process that generates, from the state npz
files saved by the Chapter-2 sweep (the Δτ=2π/N runs), three grid
figures reading out the per-edge complex waves z\_e ∈ C on the complex
plane --- (1) step 0, (2) the final step (step 500), and (3) a zoom into
the largest angular cluster at the final step. The design intent of the
three figures (their division of roles corroborating the start, end, and
end-interior of the inflation figure from the configuration side, the
caution that ``the complex plane of the figures ≠ the parent plane Π'',
and the bridging identity) is given in \textbf{Chapter 2, §2.6 (Eq.
25)}; this chapter fixes the implementation, execution, and
observations. The data are read-only; the dynamics and stored data are
never touched.

\subsection{2. Theoretical Background}\label{theoretical-background}

The mathematics of states and measurement follows Chapter 2 (Eqs.
22--25). The figure-specific conventions of this chapter are defined as
Eqs. 26--28.

\textbf{(Eq. 26) Duplicate counting (notation of degeneracy)} --- The
set of plotted complex values \{z\_e\} is partitioned into equivalence
classes by \textbf{rounding the real and imaginary parts to the 12th
decimal place}; classes with c \textgreater{} 1 members are annotated
``xc'' at their position:

\begin{verbatim}
class(w) = ( round(Re w, 12), round(Im w, 12) )
\end{verbatim}

Rounding at 12 digits is coarser than the effective double precision
(\textasciitilde16 digits): it identifies differences at the level of
rounding error while keeping physically significant separations
(relative 10\^{}-4 or larger in this series) distinct. In the zoom
figure the rounding is tightened to the 15th decimal (Eq. 28), so that
only machine-precision coincidences are counted as ``xn''.

\textbf{(Eq. 27) Panel scaling convention} --- The axis range of each
panel is the square {[}−1.15r, +1.15r{]} (equal aspect) with r = max\_e
\textbar z\_e\textbar{} the overall amplitude of that panel, and
\textbf{the tick values display the actual values} (no relabeling by
normalized values). The line segments from the origin to each point are
guides for reading the argument and modulus of each wave simultaneously.

\textbf{(Eq. 28) Extraction of the largest angular cluster (algorithm of
the zoom figure)} --- The final-step state, in coordinates normalized by
the overall amplitude, is grouped by rounding at \textbf{coarse
resolution 1/100}; the group with the most members is the zoom target:

\begin{verbatim}
key(w) = ( round(Re w / amp, 2), round(Im w / amp, 2) ),   amp = max_e |z_e|
C* = argmax_{key} |{ w : key(w) = key }|
center c = mean(C*),  spread = max_{w∈C*} |w − c|,  window = 1.4 × spread
\end{verbatim}

The panel title prints \textbar C*\textbar, spread, and amp. The spread
quantifies whether the cluster is an exact single point (exact
degeneracy) or a bundle of finite width (the zoom row of the table in
Chapter 2, §2.6).

\subsection{3. Implementation Method}\label{implementation-method}

\begin{itemize}
\tightlist
\item
  Plotting program
  \texttt{plot\_complex\_plane\_N3\_N40\_stage123\_v1.py} (bundled in
  this package; read-only). The drawing style extends the existing
  plotting programs of this series (the grid style of
  \texttt{complex\_plane\_readout\_step0\_step2000\_20260904/plot\_complex\_plane\_step0\_step2000.py}
  and the cluster-zoom algorithm of
  \texttt{自発的分裂予備実験\_v1/N40\_state\_readout\_20260904/plot\_complex\_plane\_N40\_v1.py})
  to the 8×5 grid of N=3..40; the algorithms themselves are identical
  (the style lineage is unchanged).
\item
  The input is solely the Chapter-2 product
  \texttt{results/hm\_N\{N\}\_den\_\{N\}\_states\_500.npz}. \textbf{The
  den=N (Δτ=2π/N) run is used as the representative for each N} (step 0
  is identical across denominators, so the representative choice matters
  only for the final step; final configurations for other denominators
  can be read out likewise from the stored npz).
\item
  The output is 3 PNG files. No data are rewritten or regenerated.
\end{itemize}

\subsection{4. Detailed Design}\label{detailed-design}

\subsubsection{4.1 Overall Flow}\label{overall-flow}

\begin{verbatim}
[initialization]  fix the input folder (results/)
[loop]  draw_grid(step=0)  : draw the step-0 state on each panel for N=3..40 (Eqs. 26, 27)
        draw_grid(step=500): same for the final-step state
        zoom grid          : largest angular cluster (Eq. 28) per panel for N=3..40
[finalization]  savefig each figure; switch off the 2 unused panels
\end{verbatim}

\subsubsection{4.2 Overall Data Flow}\label{overall-data-flow}

\begin{itemize}
\tightlist
\item
  \textbf{Input}:
  \texttt{results/hm\_N\{N\}\_den\_\{N\}\_states\_500.npz} × 38
  (Chapter-2 products; read-only). Fields used: \texttt{Z} (501×M),
  \texttt{denominator}, \texttt{steps} (for consistency asserts)
\item
  \textbf{Parameters} (all constants inside the program):

  \begin{itemize}
  \tightlist
  \item
    rounding digits: 12 (Eq. 26) / 15 (inside the zoom, Eq. 28) / coarse
    resolution 2 digits (the 1/100 of Eq. 28)
  \item
    scale factor 1.15 (Eq. 27), window factor 1.4 (Eq. 28)
  \item
    grid 8×5 (38 used, 2 off), dpi=180
  \end{itemize}
\item
  \textbf{Output}:

  \begin{itemize}
  \tightlist
  \item
    \texttt{fig\_complex\_plane\_step0\_N3\_N40\_stage123.png}
  \item
    \texttt{fig\_complex\_plane\_final\_N3\_N40\_stage123.png}
  \item
    \texttt{fig\_complex\_plane\_final\_zoom\_N3\_N40\_stage123.png}
  \end{itemize}
\end{itemize}

\subsubsection{4.3 Individual Processes}\label{individual-processes}

\paragraph{4.3.1 Initialization (loading and consistency
checks)}\label{initialization-loading-and-consistency-checks}

\begin{Shaded}
\begin{Highlighting}[]
    \DecValTok{18}  \KeywordTok{def}\NormalTok{ load(N, step):}
    \DecValTok{19}\NormalTok{      d }\OperatorTok{=}\NormalTok{ np.load(os.path.join(IN, }\SpecialStringTok{f\textquotesingle{}hm\_N}\SpecialCharTok{\{}\NormalTok{N}\SpecialCharTok{\}}\SpecialStringTok{\_den\_}\SpecialCharTok{\{}\NormalTok{N}\SpecialCharTok{\}}\SpecialStringTok{\_states\_500.npz\textquotesingle{}}\NormalTok{))}
    \DecValTok{20}      \ControlFlowTok{assert} \BuiltInTok{int}\NormalTok{(d[}\StringTok{\textquotesingle{}denominator\textquotesingle{}}\NormalTok{]) }\OperatorTok{==}\NormalTok{ N }\KeywordTok{and} \BuiltInTok{int}\NormalTok{(d[}\StringTok{\textquotesingle{}steps\textquotesingle{}}\NormalTok{]) }\OperatorTok{==} \DecValTok{500}
    \DecValTok{21}      \ControlFlowTok{return}\NormalTok{ np.asarray(d[}\StringTok{\textquotesingle{}Z\textquotesingle{}}\NormalTok{][step], dtype}\OperatorTok{=}\NormalTok{np.complex128)}
\end{Highlighting}
\end{Shaded}

\begin{itemize}
\tightlist
\item
  The assert on line 20 checks that ``the intended file (den=N, 500
  steps) is being read''; on mismatch an AssertionError stops execution
  (no drawing occurs).
\end{itemize}

\paragraph{4.3.2 Loop (1): Grid Rendering (Eqs. 26,
27)}\label{loop-1-grid-rendering-eqs.-26-27}

\begin{Shaded}
\begin{Highlighting}[]
    \DecValTok{26}      \ControlFlowTok{for}\NormalTok{ k, N }\KeywordTok{in} \BuiltInTok{enumerate}\NormalTok{(}\BuiltInTok{range}\NormalTok{(}\DecValTok{3}\NormalTok{, }\DecValTok{41}\NormalTok{)):}
    \DecValTok{27}\NormalTok{          ax }\OperatorTok{=}\NormalTok{ axs[k]}
    \DecValTok{28}\NormalTok{          z }\OperatorTok{=}\NormalTok{ load(N, step)}
    \DecValTok{29}\NormalTok{          M }\OperatorTok{=}\NormalTok{ N }\OperatorTok{*}\NormalTok{ (N }\OperatorTok{{-}} \DecValTok{1}\NormalTok{) }\OperatorTok{//} \DecValTok{2}
    \DecValTok{30}          \ControlFlowTok{assert}\NormalTok{ z.size }\OperatorTok{==}\NormalTok{ M}
    \DecValTok{31}          \ControlFlowTok{for}\NormalTok{ w }\KeywordTok{in}\NormalTok{ z:}
    \DecValTok{32}\NormalTok{              ax.plot([}\FloatTok{0.0}\NormalTok{, w.real], [}\FloatTok{0.0}\NormalTok{, w.imag], color}\OperatorTok{=}\StringTok{\textquotesingle{}tab:blue\textquotesingle{}}\NormalTok{, linewidth}\OperatorTok{=}\FloatTok{0.7}\NormalTok{, alpha}\OperatorTok{=}\FloatTok{0.6}\NormalTok{)}
    \DecValTok{33}\NormalTok{          ax.plot(z.real, z.imag, }\StringTok{\textquotesingle{}o\textquotesingle{}}\NormalTok{, ms}\OperatorTok{=}\FloatTok{2.5}\NormalTok{, color}\OperatorTok{=}\StringTok{\textquotesingle{}tab:red\textquotesingle{}}\NormalTok{, alpha}\OperatorTok{=}\FloatTok{0.85}\NormalTok{, linestyle}\OperatorTok{=}\StringTok{\textquotesingle{}none\textquotesingle{}}\NormalTok{)}
    \DecValTok{34}\NormalTok{          cnt }\OperatorTok{=}\NormalTok{ Counter((}\BuiltInTok{round}\NormalTok{(}\BuiltInTok{float}\NormalTok{(w.real), }\DecValTok{12}\NormalTok{), }\BuiltInTok{round}\NormalTok{(}\BuiltInTok{float}\NormalTok{(w.imag), }\DecValTok{12}\NormalTok{)) }\ControlFlowTok{for}\NormalTok{ w }\KeywordTok{in}\NormalTok{ z)}
    \DecValTok{35}          \ControlFlowTok{for}\NormalTok{ (a, b), c }\KeywordTok{in}\NormalTok{ cnt.items():}
    \DecValTok{36}              \ControlFlowTok{if}\NormalTok{ c }\OperatorTok{\textgreater{}} \DecValTok{1}\NormalTok{:}
    \DecValTok{37}\NormalTok{                  ax.annotate(}\SpecialStringTok{f\textquotesingle{}x}\SpecialCharTok{\{}\NormalTok{c}\SpecialCharTok{\}}\SpecialStringTok{\textquotesingle{}}\NormalTok{, (a, b), textcoords}\OperatorTok{=}\StringTok{\textquotesingle{}offset points\textquotesingle{}}\NormalTok{, xytext}\OperatorTok{=}\NormalTok{(}\DecValTok{3}\NormalTok{, }\DecValTok{3}\NormalTok{),}
    \DecValTok{38}\NormalTok{                              fontsize}\OperatorTok{=}\DecValTok{5}\NormalTok{, color}\OperatorTok{=}\StringTok{\textquotesingle{}black\textquotesingle{}}\NormalTok{)}
    \DecValTok{39}\NormalTok{          r }\OperatorTok{=} \BuiltInTok{float}\NormalTok{(np.}\BuiltInTok{abs}\NormalTok{(z).}\BuiltInTok{max}\NormalTok{())}
    \DecValTok{40}\NormalTok{          lim }\OperatorTok{=}\NormalTok{ r }\OperatorTok{*} \FloatTok{1.15} \ControlFlowTok{if}\NormalTok{ r }\OperatorTok{\textgreater{}} \DecValTok{0} \ControlFlowTok{else} \FloatTok{1.0}
    \DecValTok{41}\NormalTok{          ax.set\_xlim(}\OperatorTok{{-}}\NormalTok{lim, lim)}\OperatorTok{;}\NormalTok{ ax.set\_ylim(}\OperatorTok{{-}}\NormalTok{lim, lim)}\OperatorTok{;}\NormalTok{ ax.set\_aspect(}\StringTok{\textquotesingle{}equal\textquotesingle{}}\NormalTok{)}
\end{Highlighting}
\end{Shaded}

\begin{itemize}
\tightlist
\item
  Lines 31--32: segments from the origin; line 33: points; lines 34--38:
  the duplicate annotations of Eq. 26; lines 39--41: the scaling
  convention of Eq. 27 (tick values remain actual values;
  \texttt{ticklabel\_format} on line 46 only switches the notation to
  scientific).
\end{itemize}

\paragraph{4.3.3 Loop (2): Zoom into the Largest Angular Cluster (Eq.
28)}\label{loop-2-zoom-into-the-largest-angular-cluster-eq.-28}

\begin{Shaded}
\begin{Highlighting}[]
    \DecValTok{67}\NormalTok{      z }\OperatorTok{=}\NormalTok{ load(N, }\DecValTok{500}\NormalTok{)}
    \DecValTok{68}\NormalTok{      amp }\OperatorTok{=} \BuiltInTok{float}\NormalTok{(np.}\BuiltInTok{abs}\NormalTok{(z).}\BuiltInTok{max}\NormalTok{())}
    \DecValTok{69}\NormalTok{      coarse }\OperatorTok{=}\NormalTok{ \{\}}
    \DecValTok{70}      \ControlFlowTok{for}\NormalTok{ w }\KeywordTok{in}\NormalTok{ z:}
    \DecValTok{71}\NormalTok{          key }\OperatorTok{=}\NormalTok{ (}\BuiltInTok{round}\NormalTok{(}\BuiltInTok{float}\NormalTok{(w.real) }\OperatorTok{/}\NormalTok{ amp, }\DecValTok{2}\NormalTok{), }\BuiltInTok{round}\NormalTok{(}\BuiltInTok{float}\NormalTok{(w.imag) }\OperatorTok{/}\NormalTok{ amp, }\DecValTok{2}\NormalTok{))}
    \DecValTok{72}\NormalTok{          coarse.setdefault(key, []).append(w)}
    \DecValTok{73}\NormalTok{      mem }\OperatorTok{=} \BuiltInTok{max}\NormalTok{(coarse.values(), key}\OperatorTok{=}\BuiltInTok{len}\NormalTok{)}
    \DecValTok{74}\NormalTok{      zz }\OperatorTok{=}\NormalTok{ np.array(mem)}
    \DecValTok{75}\NormalTok{      c }\OperatorTok{=}\NormalTok{ zz.mean()}
    \DecValTok{76}\NormalTok{      dev }\OperatorTok{=}\NormalTok{ np.}\BuiltInTok{abs}\NormalTok{(zz }\OperatorTok{{-}}\NormalTok{ c)}
    \DecValTok{77}\NormalTok{      spread }\OperatorTok{=} \BuiltInTok{float}\NormalTok{(dev.}\BuiltInTok{max}\NormalTok{())}
    \DecValTok{78}\NormalTok{      win }\OperatorTok{=}\NormalTok{ spread }\OperatorTok{*} \FloatTok{1.4} \ControlFlowTok{if}\NormalTok{ spread }\OperatorTok{\textgreater{}} \DecValTok{0} \ControlFlowTok{else}\NormalTok{ amp }\OperatorTok{*} \FloatTok{1e{-}12}
\NormalTok{...}
    \DecValTok{90}\NormalTok{      ax.set\_title(}\SpecialStringTok{f\textquotesingle{}N=}\SpecialCharTok{\{}\NormalTok{N}\SpecialCharTok{\}}\SpecialStringTok{: }\SpecialCharTok{\{}\BuiltInTok{len}\NormalTok{(zz)}\SpecialCharTok{\}}\SpecialStringTok{ waves, dev=}\SpecialCharTok{\{}\NormalTok{spread}\SpecialCharTok{:.2e\}}\SpecialStringTok{ (|z|max=}\SpecialCharTok{\{}\NormalTok{amp}\SpecialCharTok{:.2e\}}\SpecialStringTok{)\textquotesingle{}}\NormalTok{, fontsize}\OperatorTok{=}\DecValTok{7}\NormalTok{)}
\end{Highlighting}
\end{Shaded}

\begin{itemize}
\tightlist
\item
  Lines 70--73: the grouping and largest-cluster selection of Eq. 28;
  lines 75--78: center, spread, window; line 90: printing member count,
  spread, and amp in the panel title (this printed output is the source
  of the observed values in §6).
\end{itemize}

\paragraph{4.3.4 Finalization and Exceptional Behavior (elements not
formulated as
equations)}\label{finalization-and-exceptional-behavior-elements-not-formulated-as-equations}

\begin{itemize}
\tightlist
\item
  Line 40, \texttt{lim\ =\ r*1.15\ if\ r\ \textgreater{}\ 0\ else\ 1.0}:
  fallback for the degenerate case of all points at the origin (does not
  occur in these data).
\item
  Line 78,
  \texttt{win\ =\ spread*1.4\ if\ spread\ \textgreater{}\ 0\ else\ amp*1e-12}:
  window fallback when the cluster has a single member (spread=0; occurs
  at N=3, where the largest cluster is a single wave).
\item
  Lines 50--51 and 91--92: of the 8×5=40 panels, the unused 39th and
  40th are set to
  \texttt{axis(\textquotesingle{}off\textquotesingle{})}.
\item
  The only exception paths are the asserts (lines 20, 30); there are no
  other branches.
\end{itemize}

\subsection{5. Execution Results}\label{execution-results}

\subsubsection{5.1 Reproduction Commands}\label{reproduction-commands}

\begin{Shaded}
\begin{Highlighting}[]
\BuiltInTok{cd}\NormalTok{ N3\_N40\_stage123\_sweep\_20260905}
\ExtensionTok{python3}\NormalTok{ plot\_complex\_plane\_N3\_N40\_stage123\_v1.py}
\CommentTok{\# or as the final step of ./run\_all.sh}
\end{Highlighting}
\end{Shaded}

\subsubsection{5.2 Execution Environment}\label{execution-environment}

\begin{itemize}
\tightlist
\item
  Python 3.9.6 (\texttt{.venv/bin/python3}), numpy 2.0.2, matplotlib
  (Agg rendering)
\item
  macOS 26.3.1 (arm64)
\end{itemize}

\subsubsection{5.3 Execution Time}\label{execution-time}

\textbf{Under about 1 minute} for the three figures in total (dominated
by 38×2 npz loads and the rendering of up to 780 segments × 38 panels).

\subsubsection{5.4 Verification Gates}\label{verification-gates}

{\def\LTcaptype{none} % do not increment counter
\begin{longtable}[]{@{}
  >{\raggedright\arraybackslash}p{(\linewidth - 6\tabcolsep) * \real{0.2500}}
  >{\raggedright\arraybackslash}p{(\linewidth - 6\tabcolsep) * \real{0.2500}}
  >{\raggedright\arraybackslash}p{(\linewidth - 6\tabcolsep) * \real{0.2500}}
  >{\raggedright\arraybackslash}p{(\linewidth - 6\tabcolsep) * \real{0.2500}}@{}}
\toprule\noalign{}
\begin{minipage}[b]{\linewidth}\raggedright
Gate
\end{minipage} & \begin{minipage}[b]{\linewidth}\raggedright
Pass condition
\end{minipage} & \begin{minipage}[b]{\linewidth}\raggedright
Measured
\end{minipage} & \begin{minipage}[b]{\linewidth}\raggedright
Verdict
\end{minipage} \\
\midrule\noalign{}
\endhead
\bottomrule\noalign{}
\endlastfoot
G1: input consistency & asserts denominator==N, steps==500, z.size==M
pass at every load & passed for all 38 N × 3 figures (no AssertionError)
& \textbf{PASS} \\
G2: completion & \texttt{ALL\ DONE} printed, 3 PNGs generated &
confirmed & \textbf{PASS} \\
\end{longtable}
}

(The bit-level lineage of the input npz themselves is guaranteed by
Chapter 2 G1 --- \texttt{Z{[}0{]}} of all 228 runs bit-identical to the
static parents.)

\subsubsection{5.5 Data}\label{data}

{\def\LTcaptype{none} % do not increment counter
\begin{longtable}[]{@{}
  >{\raggedright\arraybackslash}p{(\linewidth - 2\tabcolsep) * \real{0.5000}}
  >{\raggedright\arraybackslash}p{(\linewidth - 2\tabcolsep) * \real{0.5000}}@{}}
\toprule\noalign{}
\begin{minipage}[b]{\linewidth}\raggedright
Item
\end{minipage} & \begin{minipage}[b]{\linewidth}\raggedright
Content
\end{minipage} \\
\midrule\noalign{}
\endhead
\bottomrule\noalign{}
\endlastfoot
Input & \texttt{results/hm\_N\{N\}\_den\_\{N\}\_states\_500.npz} × 38
(Chapter-2 products; read-only) \\
Output location & package root \\
step-0 figure &
\texttt{fig\_complex\_plane\_step0\_N3\_N40\_stage123.png} (632,949
bytes) \\
final-step figure &
\texttt{fig\_complex\_plane\_final\_N3\_N40\_stage123.png} (4,582,367
bytes) \\
zoom figure &
\texttt{fig\_complex\_plane\_final\_zoom\_N3\_N40\_stage123.png}
(494,906 bytes) \\
SHA256 & the bundled \texttt{SHA256SUMS.txt} is canonical \\
\end{longtable}
}

\subsubsection{5.6 Figures}\label{figures}

The three figures above (8×5 grids; each panel equal aspect,
actual-value ticks). How to read them, and their correspondence with the
inflation figure, follow the table in Chapter 2, §2.6.

\subsection{6. Execution Analysis (objective report and observations
only; the
numbers}\label{execution-analysis-objective-report-and-observations-only-the-numbers}

originate from the in-figure printing = the output of this program)

\begin{enumerate}
\def\labelenumi{\arabic{enumi}.}
\tightlist
\item
  \textbf{Step-0 figure}: every panel for N=3..40 shows the star shape
  of two antipodal pairs (4 bundles). The bundles have a radial spread
  of amplitudes. The duplicate annotations of Eq. 26 appear only for
  some small N (x4 at N=4; x2--x3 for N=5--9, etc.); exact duplicates
  vanish as N grows.
\item
  \textbf{Final-step figure}: in every panel the star has disappeared,
  replaced by a nearly equimodular ring-like arrangement (sparse spokes
  for small N). The Eq. 26 (12-digit) duplicate annotations do not
  appear except for a few small-N cases (x2 at N=4, 5).
\item
  \textbf{Zoom figure}: the largest angular cluster has 1--10 members
  (N=3: 1 wave; N=39, 40: 10 waves); the spread is of order
  10\textsuperscript{-7--10}-4 for N≥6 (e.g., at N=40, dev=2.04e-04 with
  \textbar z\textbar max=3.64e-02, relative
  \textasciitilde5.6×10\^{}-3). Only the two-wave clusters of N=4, 5
  coincide at machine precision (dev 10\textsuperscript{-16--10}-10).
\item
  Summary of observations (facts only): for all N≥6 the final-step
  angular clusters are bundles of finite width, not condensation onto
  exactly identical complex values. As per the division of roles in
  Chapter 2 §2.6, this records the fine organization of the final state
  that H⊥/H does not measure.
\end{enumerate}

\begin{center}\rule{0.5\linewidth}{0.5pt}\end{center}

(End of Chapter 3. The overview paper is separate.)

\end{document}
